\section{Introduction}
Perturbation atlases measure selected interventions in selected cellular contexts, leaving most intervention--context combinations unobserved. Genome-scale Perturb-seq and multimodal perturbation screens provide unusually rich settings for studying this missingness \citep{replogle2022,frangieh2021}, while harmonized releases make computational replay practical \citep{scperturb2024}. We consider a particular information interface: the response to each query intervention has been measured in a supplied source context, but its recipient response remains unavailable. A limited recipient budget can purchase responses to other calibration interventions and additional controls. The choice is not only which interventions to measure, but how deeply to measure each one and how much to spend on the shared control reference.

These choices couple prediction, measurement precision and transfer assumptions. Another calibration intervention can improve coverage of source-response directions; more cells reduce sampling noise; more shared controls improve all control-relative contrasts together. However, a well-conditioned fit in estimated source coordinates can coexist with misspecified transfer or errors in those coordinates. A design criterion that reduces a surrogate variance need not improve a strongly regularized predictor's squared error. Distinguishing these possibilities requires matching information, prediction backbones and actual purchases, rather than comparing method names with different access to recipient labels.

The constituent ideas have substantial precedent. IterPert uses multimodal and cross-context priors for sequential perturbation acquisition \citep{iterpert2024}. PerturbPlan studies prospective power/cost trade-offs involving cell and read depth \citep{perturbplan2026}; its association-testing objective differs from sealed-query response prediction. PerturbMap transports the same intervention's source response through calibrated cross-context experts \citep{perturbmap2026}. Shared-response references, gene-coordinate permutation and representation-coverage diagnostics have precedent in AmbiModBench \citep{ambimodbench2026}; we use related checks to qualify a different, costed acquisition estimand. Transductive experimental design separates measurement and query sets \citep{fiez2019}, and kernel posterior variance is classical \citep{gpml2006}. We claim neither priority for these components nor reproduction of the cited methods' published performance. Our comparators are specified implementations under one common replay interface; the transfer reconstruction's differences are disclosed.

Our empirical question is narrower: does adding a fixed source-coordinate radial-basis-function (RBF) discrepancy-variance term to a covariance-aware greedy worst-query design improve a shared paid-validation ridge predictor? The added term is a concrete hypothesis, not an estimator of realized structural bias or a fitted discrepancy component of the prediction model. We test it in known-truth simulations, a separate development context, and a frozen bidirectional released-data replay. Post-primary analyses and a separately declared split/rank/budget robustness study examine limitations without replacing the original comparison.

The contribution is a precisely specified, reproducible evaluation with negative and failure evidence. We expose the roles of shrinkage, shared controls, estimated support and evaluation noise; preserve complete crossed baselines and paid ledgers; and distinguish absolute-effect error from ordinal ranking. The real-data evidence concerns a count-selected cohort and dependent replay tasks, not independent biological replications or prospective laboratory savings. In particular, the original real pilot already spans the estimated rank-four representation: those runs principally test precision allocation conditional on estimated support, rather than an identification transition.
